% Generated draft: requires mathematical and source review.
\documentclass[11pt]{article}
\usepackage[margin=25mm]{geometry}
\usepackage{fontspec}
\setmainfont{Latin Modern Roman}
\usepackage{amsmath,amssymb,mathtools,unicode-math}
\setmathfont{Latin Modern Math}
\usepackage{enumitem,needspace,xurl,hyperref}
\setlength{\parindent}{0pt}
\setlength{\parskip}{5pt}
\setlength{\emergencystretch}{3em}
\newcommand{\N}{\mathbb{N}}
\newcommand{\Z}{\mathbb{Z}}
\newcommand{\Q}{\mathbb{Q}}
\newcommand{\R}{\mathbb{R}}
\newcommand{\C}{\mathbb{C}}
\newcommand{\F}{\mathbb{F}}
\newcommand{\A}{\mathbb{A}}
\newcommand{\PP}{\mathbb{P}}
\newcommand{\E}{\mathbb{E}}
\newcommand{\eps}{\varepsilon}
\newcommand{\dd}{\,\mathrm{d}}
\begin{document}
\section*{UnsolvedMath ID: 800802}
\textbf{Source problem code:} catalogue entry
\subsection{Definitions and mathematical statement}
Let \(M\) denote the fixed domain occupied by the fluid, containing sources and sinks. Here sources and sinks mean locations or regions where fluid enters or leaves the system; the supplied record does not specify their mathematical representation, location, or prescribed fluxes. The term \emph{ideal fluid} is retained from the source without adding assumptions about compressibility, density, or governing equations.

\textbf{Formulation problem.} Find a variational formulation describing the motion of an ideal fluid in \(M\) with sources and sinks. Such a formulation should specify its configuration variables, admissible histories, variational functional, and allowed variations, so that its stationarity conditions describe the intended fluid motion. Preferably, find a geodesic formulation, meaning a description of the motion by geodesics of a suitable geometric configuration space.

As part of determining the correct variational problem, determine whether and how exterior forces and/or the fluid's ``memory'' must be taken into account. The source does not define memory or specify whether exterior forces are prescribed, induced by the sources and sinks, or otherwise modeled. Their inclusion is a question to be resolved, not an assumption already imposed.

\textbf{Unresolved qualifications.} The supplied statement does not specify the dimension or geometry of \(M\), boundary conditions, regularity, source and sink data, or the precise equations whose motion is to be recovered. Consequently, it specifies a research task rather than a uniquely determined variational existence assertion. A precise choice of these data requires additional source information or an explicitly qualified interpretation.
\subsection{Short English statement}
Find a variational description, preferably a geodesic one, for an ideal fluid moving in a fixed domain with sources and sinks. Determine whether exterior forces and/or fluid ``memory'' belong in that description. The supplied record leaves the fluid model, source and sink mechanism, and meaning of memory unspecified.
\subsection{Sources}
\begin{itemize}
\item Local UnsolvedMath registry, permanent ID \(800802\), source problem code \texttt{catalogue entry}, title \emph{Geometric Hydrodynamics in Open Problems --- Problem 1}. The supplied record attributes the problem to Khesin, Misiolek, and Modin, with proposed year 2022, and labels it \emph{open}; this is a supplied registry status, not an independently verified mathematical or literature assessment.
\item Khesin, Misiolek, and Modin, \emph{Geometric Hydrodynamics in Open Problems} (2022), Problem 1. Supplied source URL: \url{https://arxiv.org/abs/2205.01143}. The provenance reports access on 2026-10-16 and extraction by \texttt{arxiv\_latex}; this access information is supplied provenance and does not indicate independent source consultation for this draft.
\item Supplied upstream record ID \(800802\), revision \texttt{d371ee917ac9647fbb06cd6fc529deec6f3dda82}; registry creation and update timestamps are both \texttt{2026-10-16T14:00:00Z}.
\end{itemize}
\end{document}
